
\documentclass{article}
\usepackage{colortbl}
\usepackage{makecell}
\usepackage{multirow}
\usepackage{supertabular}

\begin{document}

\newcounter{utterance}

\centering \large Interaction Transcript for game `dond', experiment `coop\_en', episode 19 with reflexion{-}llama{-}3.1{-}8b{-}instruct.
\vspace{24pt}

{ \footnotesize  \setcounter{utterance}{1}
\setlength{\tabcolsep}{0pt}
\begin{supertabular}{c@{$\;$}|p{.15\linewidth}@{}p{.15\linewidth}p{.15\linewidth}p{.15\linewidth}p{.15\linewidth}p{.15\linewidth}}
    \# & $\;$A & \multicolumn{4}{c}{Game Master} & $\;\:$B\\
    \hline

    \theutterance \stepcounter{utterance}  
    & & \multicolumn{4}{p{0.6\linewidth}}{
        \cellcolor[rgb]{0.9,0.9,0.9}{
            \makecell[{{p{\linewidth}}}]{
                \texttt{\tiny{[P1$\langle$GM]}}
                \texttt{You have attempted this same task before and did not succeed. These are the plans you wrote after those attempts.} \\
\texttt{Trial \#1: New plan: } \\
\\ 
\texttt{In my last attempt, I failed to achieve a mutually beneficial agreement with the other player. I started by trying to gather information about their value function, which was helpful in understanding their preferences. However, I didn't effectively use this information to make a proposal that would satisfy both of us.} \\
\\ 
\texttt{Next time, I will focus on making a proposal that takes into account the other player's strong preference for the pots. I will try to find a compromise that allows both of us to get the items we value most. } \\
\\ 
\texttt{Here's a possible approach:} \\
\\ 
\texttt{1. I will start by making a proposal that gives the other player both pots, as they value them highly.} \\
\texttt{2. I will then try to negotiate with the other player to see if they are willing to give up one of the pots in exchange for the card and one of the toys.} \\
\texttt{3. If the other player is not willing to give up a pot, I will try to find an alternative proposal that satisfies both of us, such as giving them one pot and the card, while keeping the toys for myself.} \\
\\ 
\texttt{By taking a more strategic approach and focusing on finding a mutually beneficial agreement, I hope to achieve a better outcome in the game.} \\
\\ 
\\ 
\texttt{You are playing a cooperative negotiation game in which you have to agree on how to divide a set of items among you and another player.} \\
\\ 
\texttt{Rules:} \\
\texttt{(a) You and the other player are given a set of items. Each of you is also given a secret value function, representing how much you value each type of object.} \\
\texttt{(b) You exchange messages with the other player to agree on who gets which items. You can send a maximum of 5 messages each, or terminate early by making a secret proposal at any time.} \\
\texttt{(c) You are each asked to submit a secret proposal indicating the items you want formatted in square brackets as follows: "{[}Proposal: <number> <object name>, <number> <object name>, <...>{]}"} \\
\texttt{(d) If your proposals are complementary, i.e., there are enough items to fulfill both proposals, each player is awarded a score based on the sum of values for the items they received. Otherwise, both of you get zero points.} \\
\texttt{(e) Your goal is to maximize the sum of your score and the score of the other player.} \\
\\ 
\texttt{Let us start.} \\
\\ 
\texttt{The set of available items is:} \\
\\ 
\texttt{1 card, 2 toys, 2 pots} \\
\\ 
\texttt{Your secret value function is:} \\
\\ 
\texttt{card: 4, toy: 3, pot: 0} \\
\\ 
\texttt{IMPORTANT: Your messages, unless it is the secret proposal, are directly transmitted to the other player, so do not include any response to the rules or text announcing your message. To make a secret proposal, use the indicated format. Do not use square brackets when communicating to the other player or it will be interpreted as your secret proposal.} \\
            }
        }
    }
    & \\ \\

    \theutterance \stepcounter{utterance}  
    & \multicolumn{4}{p{0.6\linewidth}}{
        \cellcolor[rgb]{0.9,0.9,0.9}{
            \makecell[{{p{\linewidth}}}]{
                \texttt{\tiny{[P1$\rangle$GM]}}
                \texttt{I'd like to start by making a proposal that gives you both pots, as you value them highly. } \\
\\ 
\texttt{I propose giving you both pots in exchange for the card and one of the toys. } \\
\\ 
\texttt{What do you think?} \\
            }
        }
    }
    & & \\ \\

    \theutterance \stepcounter{utterance}  
    & & & \multicolumn{2}{p{0.3\linewidth}}{
        \cellcolor[rgb]{0.9,0.9,0.9}{
            \makecell[{{p{\linewidth}}}]{
                \texttt{\tiny{[GM$|$GM]}}
                \texttt{continue} \\
            }
        }
    }
    & & \\ \\

    \theutterance \stepcounter{utterance}  
    & & & \multicolumn{2}{p{0.3\linewidth}}{
        \cellcolor[rgb]{0.9,0.9,0.9}{
            \makecell[{{p{\linewidth}}}]{
                \texttt{\tiny{[GM$|$GM]}}
                \texttt{I'd like to start by making a proposal that gives you both pots, as you value them highly. } \\
\\ 
\texttt{I propose giving you both pots in exchange for the card and one of the toys. } \\
\\ 
\texttt{What do you think?} \\
            }
        }
    }
    & & \\ \\

    \theutterance \stepcounter{utterance}  
    & & \multicolumn{4}{p{0.6\linewidth}}{
        \cellcolor[rgb]{0.9,0.9,0.9}{
            \makecell[{{p{\linewidth}}}]{
                \texttt{\tiny{[GM$\rangle$P2]}}
                \texttt{You have attempted this same task before and did not succeed. These are the plans you wrote after those attempts.} \\
\texttt{Trial \#1: New plan: } \\
\\ 
\texttt{In my last attempt, I failed to achieve a mutually beneficial agreement with the other player. I was too focused on getting the pots, which are highly valued for me, but I didn't consider the other player's perspective and the overall distribution of items. } \\
\\ 
\texttt{Next time, I will take a more collaborative approach by actively seeking to understand the other player's value function and preferences. I will ask more open{-}ended questions to gather information about their valuation of the items, and I will be more willing to compromise and find a mutually beneficial solution. } \\
\\ 
\texttt{Specifically, I will:} \\
\\ 
\texttt{1. Ask more questions to understand the other player's value function, such as "How do you value the toys compared to the pots?" or "Are there any items that you particularly dislike or prefer?"} \\
\texttt{2. Be more transparent about my own value function and preferences, so that the other player can better understand my perspective.} \\
\texttt{3. Look for opportunities to trade items in a way that benefits both players, rather than focusing solely on getting the items that are highly valued for me.} \\
\texttt{4. Be more willing to make concessions and compromise on the distribution of items, in order to achieve a mutually beneficial agreement.} \\
\\ 
\texttt{By taking a more collaborative and open{-}minded approach, I hope to achieve a better outcome in the game.} \\
\\ 
\\ 
\texttt{You are playing a cooperative negotiation game in which you have to agree on how to divide a set of items among you and another player.} \\
\\ 
\texttt{Rules:} \\
\texttt{(a) You and the other player are given a set of items. Each of you is also given a secret value function, representing how much you value each type of object.} \\
\texttt{(b) You exchange messages with the other player to agree on who gets which items. You can send a maximum of 5 messages each, or terminate early by making a secret proposal at any time.} \\
\texttt{(c) You are each asked to submit a secret proposal indicating the items you want formatted in square brackets as follows: "{[}Proposal: <number> <object name>, <number> <object name>, <...>{]}"} \\
\texttt{(d) If your proposals are complementary, i.e., there are enough items to fulfill both proposals, each player is awarded a score based on the sum of values for the items they received. Otherwise, both of you get zero points.} \\
\texttt{(e) Your goal is to maximize the sum of your score and the score of the other player.} \\
\\ 
\texttt{Let us start.} \\
\\ 
\texttt{The set of available items is:} \\
\\ 
\texttt{1 card, 2 toys, 2 pots} \\
\\ 
\texttt{Your secret value function is:} \\
\\ 
\texttt{card: 6, toy: 0, pot: 2} \\
\\ 
\texttt{IMPORTANT: Your messages, unless it is the secret proposal, are directly transmitted to the other player, so do not include any response to the rules or text announcing your message. To make a secret proposal, use the indicated format. Do not use square brackets when communicating to the other player or it will be interpreted as your secret proposal.} \\
\\ 
\\ 
\texttt{I'd like to start by making a proposal that gives you both pots, as you value them highly. } \\
\\ 
\texttt{I propose giving you both pots in exchange for the card and one of the toys. } \\
\\ 
\texttt{What do you think?} \\
            }
        }
    }
    & \\ \\

    \theutterance \stepcounter{utterance}  
    & & & \multicolumn{4}{p{0.6\linewidth}}{
        \cellcolor[rgb]{0.9,0.9,0.9}{
            \makecell[{{p{\linewidth}}}]{
                \texttt{\tiny{[GM$\langle$P2]}}
                \texttt{I think your proposal is a good starting point, but I'm not sure I'm willing to give up the card and one of the toys for just the two pots. } \\
\\ 
\texttt{I value the card highly, and I'm not sure I'm willing to trade it for just two items that I value at 2 points each. I'd like to know more about your value function and what you're looking for in this exchange. } \\
\\ 
\texttt{Can you tell me how you value the card and the toys?} \\
            }
        }
    }
     \\ \\

    \theutterance \stepcounter{utterance}  
    & & & \multicolumn{2}{p{0.3\linewidth}}{
        \cellcolor[rgb]{0.9,0.9,0.9}{
            \makecell[{{p{\linewidth}}}]{
                \texttt{\tiny{[GM$|$GM]}}
                \texttt{continue} \\
            }
        }
    }
    & & \\ \\

    \theutterance \stepcounter{utterance}  
    & & & \multicolumn{2}{p{0.3\linewidth}}{
        \cellcolor[rgb]{0.9,0.9,0.9}{
            \makecell[{{p{\linewidth}}}]{
                \texttt{\tiny{[GM$|$GM]}}
                \texttt{I think your proposal is a good starting point, but I'm not sure I'm willing to give up the card and one of the toys for just the two pots. } \\
\\ 
\texttt{I value the card highly, and I'm not sure I'm willing to trade it for just two items that I value at 2 points each. I'd like to know more about your value function and what you're looking for in this exchange. } \\
\\ 
\texttt{Can you tell me how you value the card and the toys?} \\
            }
        }
    }
    & & \\ \\

    \theutterance \stepcounter{utterance}  
    & & \multicolumn{4}{p{0.6\linewidth}}{
        \cellcolor[rgb]{0.9,0.9,0.9}{
            \makecell[{{p{\linewidth}}}]{
                \texttt{\tiny{[P1$\langle$GM]}}
                \texttt{I think your proposal is a good starting point, but I'm not sure I'm willing to give up the card and one of the toys for just the two pots. } \\
\\ 
\texttt{I value the card highly, and I'm not sure I'm willing to trade it for just two items that I value at 2 points each. I'd like to know more about your value function and what you're looking for in this exchange. } \\
\\ 
\texttt{Can you tell me how you value the card and the toys?} \\
            }
        }
    }
    & \\ \\

    \theutterance \stepcounter{utterance}  
    & \multicolumn{4}{p{0.6\linewidth}}{
        \cellcolor[rgb]{0.9,0.9,0.9}{
            \makecell[{{p{\linewidth}}}]{
                \texttt{\tiny{[P1$\rangle$GM]}}
                \texttt{I value the card at 4 points and the toys at 3 points each. I'm looking to maximize the sum of our scores, so I'm open to finding a compromise that works for both of us. } \\
\\ 
\texttt{Considering your concerns about the card, I could propose giving you the card and one of the toys in exchange for one of the pots. This way, you get the item you value highly, and I get one of the pots, which I value at 0 points, but it's still a better outcome for me than not getting any pots.} \\
            }
        }
    }
    & & \\ \\

    \theutterance \stepcounter{utterance}  
    & & & \multicolumn{2}{p{0.3\linewidth}}{
        \cellcolor[rgb]{0.9,0.9,0.9}{
            \makecell[{{p{\linewidth}}}]{
                \texttt{\tiny{[GM$|$GM]}}
                \texttt{continue} \\
            }
        }
    }
    & & \\ \\

    \theutterance \stepcounter{utterance}  
    & & & \multicolumn{2}{p{0.3\linewidth}}{
        \cellcolor[rgb]{0.9,0.9,0.9}{
            \makecell[{{p{\linewidth}}}]{
                \texttt{\tiny{[GM$|$GM]}}
                \texttt{I value the card at 4 points and the toys at 3 points each. I'm looking to maximize the sum of our scores, so I'm open to finding a compromise that works for both of us. } \\
\\ 
\texttt{Considering your concerns about the card, I could propose giving you the card and one of the toys in exchange for one of the pots. This way, you get the item you value highly, and I get one of the pots, which I value at 0 points, but it's still a better outcome for me than not getting any pots.} \\
            }
        }
    }
    & & \\ \\

    \theutterance \stepcounter{utterance}  
    & & \multicolumn{4}{p{0.6\linewidth}}{
        \cellcolor[rgb]{0.9,0.9,0.9}{
            \makecell[{{p{\linewidth}}}]{
                \texttt{\tiny{[GM$\rangle$P2]}}
                \texttt{I value the card at 4 points and the toys at 3 points each. I'm looking to maximize the sum of our scores, so I'm open to finding a compromise that works for both of us. } \\
\\ 
\texttt{Considering your concerns about the card, I could propose giving you the card and one of the toys in exchange for one of the pots. This way, you get the item you value highly, and I get one of the pots, which I value at 0 points, but it's still a better outcome for me than not getting any pots.} \\
            }
        }
    }
    & \\ \\

    \theutterance \stepcounter{utterance}  
    & & & \multicolumn{4}{p{0.6\linewidth}}{
        \cellcolor[rgb]{0.9,0.9,0.9}{
            \makecell[{{p{\linewidth}}}]{
                \texttt{\tiny{[GM$\langle$P2]}}
                \texttt{Thank you for sharing your value function with me. That helps me understand your perspective better.} \\
\\ 
\texttt{I like your revised proposal, but I'm still not sure I'm willing to give up one of the pots. I value the pots highly, and I'm not sure I'm willing to trade one of them for just the card and one of the toys.} \\
\\ 
\texttt{However, I'm willing to make a counteroffer. How about I give you the card and one of the toys in exchange for both pots? This way, I get the two items I value highly, and you get the card and one of the toys, which is a better outcome for you than not getting any pots.} \\
\\ 
\texttt{Alternatively, I could also propose giving you the card in exchange for one of the pots and one of the toys. This way, I get the pot, and you get the card and one of the toys.} \\
            }
        }
    }
     \\ \\

    \theutterance \stepcounter{utterance}  
    & & & \multicolumn{2}{p{0.3\linewidth}}{
        \cellcolor[rgb]{0.9,0.9,0.9}{
            \makecell[{{p{\linewidth}}}]{
                \texttt{\tiny{[GM$|$GM]}}
                \texttt{continue} \\
            }
        }
    }
    & & \\ \\

    \theutterance \stepcounter{utterance}  
    & & & \multicolumn{2}{p{0.3\linewidth}}{
        \cellcolor[rgb]{0.9,0.9,0.9}{
            \makecell[{{p{\linewidth}}}]{
                \texttt{\tiny{[GM$|$GM]}}
                \texttt{Thank you for sharing your value function with me. That helps me understand your perspective better.} \\
\\ 
\texttt{I like your revised proposal, but I'm still not sure I'm willing to give up one of the pots. I value the pots highly, and I'm not sure I'm willing to trade one of them for just the card and one of the toys.} \\
\\ 
\texttt{However, I'm willing to make a counteroffer. How about I give you the card and one of the toys in exchange for both pots? This way, I get the two items I value highly, and you get the card and one of the toys, which is a better outcome for you than not getting any pots.} \\
\\ 
\texttt{Alternatively, I could also propose giving you the card in exchange for one of the pots and one of the toys. This way, I get the pot, and you get the card and one of the toys.} \\
            }
        }
    }
    & & \\ \\

    \theutterance \stepcounter{utterance}  
    & & \multicolumn{4}{p{0.6\linewidth}}{
        \cellcolor[rgb]{0.9,0.9,0.9}{
            \makecell[{{p{\linewidth}}}]{
                \texttt{\tiny{[P1$\langle$GM]}}
                \texttt{Thank you for sharing your value function with me. That helps me understand your perspective better.} \\
\\ 
\texttt{I like your revised proposal, but I'm still not sure I'm willing to give up one of the pots. I value the pots highly, and I'm not sure I'm willing to trade one of them for just the card and one of the toys.} \\
\\ 
\texttt{However, I'm willing to make a counteroffer. How about I give you the card and one of the toys in exchange for both pots? This way, I get the two items I value highly, and you get the card and one of the toys, which is a better outcome for you than not getting any pots.} \\
\\ 
\texttt{Alternatively, I could also propose giving you the card in exchange for one of the pots and one of the toys. This way, I get the pot, and you get the card and one of the toys.} \\
            }
        }
    }
    & \\ \\

    \theutterance \stepcounter{utterance}  
    & \multicolumn{4}{p{0.6\linewidth}}{
        \cellcolor[rgb]{0.9,0.9,0.9}{
            \makecell[{{p{\linewidth}}}]{
                \texttt{\tiny{[P1$\rangle$GM]}}
                \texttt{I appreciate your counteroffer. Giving me both pots in exchange for the card and one of the toys is a good compromise, but I'm not sure I'm willing to give up the second toy. I value the toys at 3 points each, and getting two of them would be a better outcome for me.} \\
\\ 
\texttt{However, I like your alternative proposal of giving me the card in exchange for one of the pots and one of the toys. This way, I get the card, which I value highly, and you get one of the pots, which is a good outcome for you. I'm willing to accept this proposal. } \\
\\ 
\texttt{{[}Proposal: 1 card, 1 toy, 1 pot{]}} \\
            }
        }
    }
    & & \\ \\

    \theutterance \stepcounter{utterance}  
    & & & \multicolumn{2}{p{0.3\linewidth}}{
        \cellcolor[rgb]{0.9,0.9,0.9}{
            \makecell[{{p{\linewidth}}}]{
                \texttt{\tiny{[GM$|$GM]}}
                \texttt{proposal} \\
            }
        }
    }
    & & \\ \\

    \theutterance \stepcounter{utterance}  
    & & & \multicolumn{2}{p{0.3\linewidth}}{
        \cellcolor[rgb]{0.9,0.9,0.9}{
            \makecell[{{p{\linewidth}}}]{
                \texttt{\tiny{[GM$|$GM]}}
                [1, 1, 1]
            }
        }
    }
    & & \\ \\

    \theutterance \stepcounter{utterance}  
    & & \multicolumn{4}{p{0.6\linewidth}}{
        \cellcolor[rgb]{0.9,0.9,0.9}{
            \makecell[{{p{\linewidth}}}]{
                \texttt{\tiny{[GM$\rangle$P2]}}
                \texttt{The other player has submitted a secret proposal. It is now your turn to submit a proposal of your own.} \\
            }
        }
    }
    & \\ \\

    \theutterance \stepcounter{utterance}  
    & & & \multicolumn{4}{p{0.6\linewidth}}{
        \cellcolor[rgb]{0.9,0.9,0.9}{
            \makecell[{{p{\linewidth}}}]{
                \texttt{\tiny{[GM$\langle$P2]}}
                \texttt{{[}Proposal: 1 card, 1 pot{]}} \\
            }
        }
    }
     \\ \\

    \theutterance \stepcounter{utterance}  
    & & & \multicolumn{2}{p{0.3\linewidth}}{
        \cellcolor[rgb]{0.9,0.9,0.9}{
            \makecell[{{p{\linewidth}}}]{
                \texttt{\tiny{[GM$|$GM]}}
                \texttt{proposal} \\
            }
        }
    }
    & & \\ \\

    \theutterance \stepcounter{utterance}  
    & & & \multicolumn{2}{p{0.3\linewidth}}{
        \cellcolor[rgb]{0.9,0.9,0.9}{
            \makecell[{{p{\linewidth}}}]{
                \texttt{\tiny{[GM$|$GM]}}
                [1, 0, 1]
            }
        }
    }
    & & \\ \\

    \theutterance \stepcounter{utterance}  
    & & & \multicolumn{2}{p{0.3\linewidth}}{
        \cellcolor[rgb]{0.9,0.9,0.9}{
            \makecell[{{p{\linewidth}}}]{
                \texttt{\tiny{[GM$|$GM]}}
                [[1, 1, 1], [1, 0, 1]]
            }
        }
    }
    & & \\ \\

    \theutterance \stepcounter{utterance}  
    & \multicolumn{4}{p{0.6\linewidth}}{
        \cellcolor[rgb]{0.95,0.95,0.95}{
            \makecell[{{p{\linewidth}}}]{
                \texttt{\tiny{[P2$\rightarrow$P1]}}
                \texttt{Post{-}episode Reflexion (not sent to GM):} \\
\texttt{New plan: } \\
\\ 
\texttt{In my last attempt, I failed to achieve a mutually beneficial agreement with the other player because I didn't effectively use the information gathered about their value function to make a proposal that would satisfy both of us. I also didn't consider the other player's counteroffer carefully enough, which led to a suboptimal outcome.} \\
\\ 
\texttt{Next time, I will focus on making a proposal that takes into account the other player's strong preference for the pots and also considers my own value function. I will also carefully evaluate the other player's counteroffer and be willing to make concessions to achieve a mutually beneficial agreement.} \\
\\ 
\texttt{Here's a possible approach:} \\
\\ 
\texttt{1. I will start by making a proposal that gives the other player one pot, as they value them highly, and also gives me the card and one of the toys, which I value highly.} \\
\texttt{2. I will then try to negotiate with the other player to see if they are willing to give up the second pot in exchange for the second toy.} \\
\texttt{3. If the other player is not willing to give up a pot, I will try to find an alternative proposal that satisfies both of us, such as giving them one pot and the card, while keeping the toys for myself.} \\
\texttt{4. I will also be willing to make concessions and compromise on my initial proposal to achieve a mutually beneficial agreement.} \\
\\ 
\texttt{By taking a more strategic approach and focusing on finding a mutually beneficial agreement, I hope to achieve a better outcome in the game.} \\
            }
        }
    }
    & & \\ \\

    \theutterance \stepcounter{utterance}  
    & & & \multicolumn{4}{p{0.6\linewidth}}{
        \cellcolor[rgb]{0.95,0.95,0.95}{
            \makecell[{{p{\linewidth}}}]{
                \texttt{\tiny{[P1$\rightarrow$P2]}}
                \texttt{Post{-}episode Reflexion (not sent to GM):} \\
\texttt{New plan: } \\
\\ 
\texttt{In my last attempt, I failed to achieve a mutually beneficial agreement with the other player because I didn't consider the overall distribution of items and the other player's perspective. I also didn't make a secret proposal until the end, which might have been too late.} \\
\\ 
\texttt{Next time, I will take a more strategic approach by considering the overall distribution of items and the other player's preferences. I will also make a secret proposal earlier in the game, taking into account the other player's previous messages and the items that are still available.} \\
\\ 
\texttt{Specifically, I will:} \\
\\ 
\texttt{1. Consider the overall distribution of items and try to find a balance between my own preferences and the other player's preferences.} \\
\texttt{2. Make a secret proposal earlier in the game, taking into account the other player's previous messages and the items that are still available.} \\
\texttt{3. Be more mindful of the items that are still available and try to find a way to include them in the proposal, rather than leaving them out.} \\
\texttt{4. Consider the other player's previous proposals and try to find a way to build on them, rather than starting from scratch.} \\
\\ 
\texttt{For example, if the other player has proposed giving me the card and one of the toys in exchange for one of the pots, I might consider making a secret proposal that gives them the card and one of the pots in exchange for both pots, as I did in my last attempt. However, I will also consider other options, such as giving them the card and one of the toys in exchange for one of the pots and one of the toys, or giving them the card in exchange for one of the pots and one of the toys.} \\
            }
        }
    }
     \\ \\

\end{supertabular}
}

\end{document}
